% Source-derived chapter 05: Proof of the geometric Bogomolov conjecture
% Original source: geometric-bogomolov-conjecture-arbitrary-characteristics
\chapter{Proof of the geometric Bogomolov conjecture}
\label{geometric-bogomolov-conjecture-arbitrary-characteristics-chapter-05}
\source{geometric-bogomolov-conjecture-arbitrary-characteristics}

\begin{remark}[Source section]
\label{geometric-bogomolov-conjecture-arbitrary-characteristics-chapter-05-source}
\source{geometric-bogomolov-conjecture-arbitrary-characteristics}
\source{geometric-bogomolov-conjecture-arbitrary-characteristics}
This chapter reproduces the corresponding section of the retrieved
LaTeX source, including its definitions, statements, examples, and
proofs. It has not yet been translated into Lean.
\notready
\end{remark}

% The source section title was: Proof of the geometric Bogomolov conjecture
\label{sec final}
\source{geometric-bogomolov-conjecture-arbitrary-characteristics}

The goal of this section is to prove Theorem \ref{thmmain}.
Let $A/K/k$ be as in the theorem. 
By Proposition \ref{changefield}, we can assume the following condition.
\begin{itemize}
\item[(a)] the transcendence degree of $K$ over $k$ is 1. 
\end{itemize}
%Replacing $k$ by its algebraic closure in $K$, we assume that $k$ is algebraically closed in $K$.
Let $S$ be the unique (normal) projective model of $K/k$. 
We can further make the following assumptions.
\begin{itemize}
\item[(b)] $S$ is a smooth projective curve over $k$;
\item[(c)] $A$ has semi-stable reduction over $S$. 
\end{itemize}
%Both condition can be achieved by replacing $k$ by a finite extension $k'$ and replacing $K$ by some compositum $Kk'$. 
The second condition is a consequence of the semistable reduction theorem. 

Our next step is to apply Yamaki's result to make the following further assumption:
\begin{itemize}
\item[(d)] $A$ has trivial $\overline K/k$-trace and good reduction over $S$. 
\end{itemize} 
%In fact, by further replacing $K$ by a finite extension $K'$ and replacing $k$ by its algebraic closure in $K'$, we can assume that  
%the $\overline K/ k$-trace $(A^{\overline K/ k},\tr)$ can be descended to the $K/k$-trace $(A^{K/k},\tr)$. 
Denote by $A'$ the maximal abelian subvariety of $A$ that has everywhere good reduction over $S$. 
Yamaki \cite[Theorem 1.5]{Yamaki2016a} asserts that the geometric Bogomolov conjecture holds for $A$ if and only if it holds for $A'/\tr(A^{K/k}\otimes_kK)$. 
This gives the condition. 

The major contribution of this section is to prove the geometric Bogomolov conjecture
under (a)-(d), based on the Manin--Mumford conjecture in this case. 

Recall that the Manin--Mumford conjecture was proved by Raynaud \cite{Raynaud1983,Raynaud1983a} over number fields, and proved by Hrushovski \cite{Hrushovski2001} over arbitrary fields.  Hrushovski's proof relies in the model theory of difference fields. Inspired by Hrushovski's proof, Pink--R\"ossler \cite{Pink2002a,Pink2004} gave a new proof using classical algebraic geometry. We will only need the conjecture under assumption (a)-(d). 
For convenience, state it in the following case of trivial $\overline K/k$-trace.

\begin{thm}[Manin--Mumford conjecture]
\source{geometric-bogomolov-conjecture-arbitrary-characteristics}
\notready \label{MM}
\source{geometric-bogomolov-conjecture-arbitrary-characteristics}
Let $K$ be a finitely generated field over an algebraically closed field $k$. Let $A$ be an abelian variety over $K$ of trivial $\overline K/ k$-trace.
Let $X$ be a closed subvariety of $A_{\overline K}$. 
Assume that $X(\overline K)\cap A(\overline K)_\tor$ is Zariski dense in $X$. 
Then $X$ is a torsion subvariety of $A_{\overline K}$. 
\end{thm}

We also need the following well-known consequence of Zhang's fundamental inequality.

\begin{lem}
\source{geometric-bogomolov-conjecture-arbitrary-characteristics}
\notready\label{fundamental inequality} 
\source{geometric-bogomolov-conjecture-arbitrary-characteristics}
Let $K/k$ be a finitely generated field extension of transcendence degree $1$. 
Let $A$ be an abelian variety over $K$. Let $L$ be a symmetric and ample line bundle over $A$.
Let $X$ be a closed subvariety of $A_\OK$.
If $X$ contains a dense set of small points of $A/K/k$, then $\widehat h_L(X)=0$.
\end{lem}
\begin{proof}
It is a direct consequence of the fundamental inequality, which asserts that
$$
\sup_{U\subset X} \inf_{x\in U(\overline K)}  
\widehat h_L(x)
\geq   \widehat h_L(X).$$
Here the supremum goes through all open subvarieties $U$ of $X$.
The fundamental inequality was proved over number fields by Zhang \cite[Theorem 1.10]{Zhang1995}, based on his previous works \cite{Zhang1992, Zhang1995a} and as a part of this theorem of successive minima. 
It was generalized to function fields by Gubler \cite[Lemma 4.1]{Gubler2007}. 
\end{proof}



 

\subsection{Subvarieties generated by addition}

With these preparations, we are ready to prove Theorem \ref{thmmain}. 
Let $A/K/k$ and $X$ be as in the theorem. 
By the above reduction, we can assume that conditions (a)-(d) hold.
By extending $K$ and $k$ if necessary, we can assume that $X=X^*\otimes_K \overline K$ for a subvariety $X^*$ of $A$ over $K$. 
By abuse of notations, we will write $X=X^*$, viewed as a subvariety of $A$.  

In this step, we consider the sum $X$ with itself in $A$. The key assumption we are going to use is that $A$ has a trivial $\overline K/k$-trace. 

For any integer $m\geq1$, denote by $X_m$ the image of the addition morphism 
$$f_m:X^m\longrightarrow A, \quad 
(x_1,\cdots, x_m) \longmapsto x_1+ \cdots+ x_m.$$ 
Set $X_0=0$ to be the identity point of $A$.
Since $X_m$ is the image of the addition morphism
$X_{m-1}\times X\to A$, we have
$$\dim X_{m-1}\leq \dim X_m\leq \dim X_{m-1}+\dim X, \quad m\geq1.$$ 

\begin{lem}
\source{geometric-bogomolov-conjecture-arbitrary-characteristics}
\notready \label{addition}
\source{geometric-bogomolov-conjecture-arbitrary-characteristics}
There is a unique integer $r\geq 1$ such that $\dim X_{r-1}<\dim X_r$ and $X_r$ is a torsion subvariety of $A$.
\end{lem}
\begin{proof}
Fix a point $x_0\in X(K)$, which exists by replacing $K$ by a finite extension. 
Denote by 
$X_m'=X_m-mx_0$
the translation of $X_m$ by $-mx_0$ in $A$.
Note that $X_m'$ is an (irreducible) subvariety of $A$, and that the sequence 
$\{X_m'\}_m$ is increasing. 
As a consequence, there is $r\geq 1$ such that $X_m'\subsetneq X_r'$ for all $m<r$ and $X_m'= X_r'$ for all $m\geq r$.

It follows that there is a morphism $\sigma: X_r' \times X_r' \to X_{2r}' = X_r'$ induced by the addition morphism of $A$. This is sufficient to imply that $B=X_r'$ is an abelian subvariety of $A$. 
In fact, the relative dimension of $\sigma$ is equal to $\dim B$. 
By semi-continuity, the dimension of the fiber $\sigma^{-1}(0):=\{(x,y)\in B^2:x+y=0\}$ is at least $\dim B$. 
On the other hand, the first projection $\sigma^{-1}(0)\to B$ is injective on $\OK$-points, and thus bijective on $\OK$-points by comparing dimensions. 
This implies that for any $x\in B(\OK)$, $y=-x\in B(\OK)$. 
Thus we have inverse morphism $[-1]B \subset B$. 
It follows that $B$ is an abelian subvariety of $A$. 

As a consequence, $X_r=B+t$ for the abelian subvariety $B$ of $A$ and the point $t=rx_0\in A(K)$. 
Denote $C=A/B$, which is an abelian variety over $K$.
It suffices to prove that the image $t'$ of $t$ in $C(K)$ is a torsion point.
There is a surjective homomorphism $A \to B$ such that the composition $B\to A \to B$ is an isogeny. 
This induces an isogeny $A\to A'$ with $A'=B\times C$.

By assumption, $X$ contains a dense set of small points of $A/K/k$.
Since the canonical height is quadratic and positive definite up to torsion, we have 
$$
\hat h(x_1+ \cdots+ x_m)\leq  m(\hat h(x_1)+ \cdots +\hat h(x_m)), \quad x_i\in A(\OK).
$$
Then ``small points'' of $X^r$ transfer to ``small points'' of $X_r$.
As a consequence, $X_r$ contains a dense set of small points of $A/K/k$.
Then its image $X'=B+t'$ in $A'$ contains a dense set of small points of $A'/K/k$.

Take symmetric and ample line bundle $L_1$ over $B$ and $L_2$ over $C$.
Choose $L=p_1^*L_1+p_2^*L_2$, which is a symmetric and ample line bundle over $A'$, where $p_1:B\times C\to B$ and $p_2:B\times C\to B$ are the projections. 
For any point $x\in B(\OK)$, we have the heights
$$
\hat h_L(x+t')= \hat h_{L_1}(x)+\hat h_{L_2}(t') \geq \hat h_{L_2}(t') \geq 0. 
$$
This forces $\hat h_{L_2}(t')=0$ as $B+t'$ contains a dense set of small points of $A'/K/k$.

By assumption, $A$ has a trivial $\overline K/ k$-trace, so 
$B$ and $C$ have trivial $\overline K/ k$-traces.
Then $\hat h_{L_2}(t')=0$ implies that $t'\in C(K)_\tor$. 
This follows from the Northcott property (cf. \cite[Theorem 9.15]{Conrad}).
It finishes the proof.
\end{proof}



By the lemma, $X_r$ is a torsion subvariety of $A$. Replacing $X$ by a translation by a suitable torsion point, we can assume that $X_r$ is an abelian subvariety of $A$. In this process, to make $X$ to be defined over $K$, we may need to replacing $K$ by a finite extension again.
To summarize, we write this process as the following condition. 
\begin{itemize}
\item[(e)] $A'=X_r$ is an abelian subvariety of $A$ and $\dim X_{r-1}<\dim X_r$. 
\end{itemize} 



\subsection{Fibers of the addition map}

By assumption (e), the image of the summation map $f_r:X^r\to A$
is an abelian subvariety $A'$ of $A$. 
This induces a summation morphism 
$$f:X_{r-1}\times X\longrightarrow A'.$$
Denote by 
$$e=\dim X_{r-1}+\dim X-\dim A'$$ 
the relative dimension of this morphism. 
By assumption (e), we have
$$e<\dim X.$$
Denote 
$$
A'_{e+1}=\{y\in A': \dim f^{-1}(y)\geq e+1 \}.
$$
Then $A'_{e+1}$ is a closed subset of $A'$ of codimension at least 2. 
Note that $X_{r-1}\times X$ is a natural subvariety of the abelian variety $A^2=A\times A$. 
The goal of this step is to prove the following result.

\begin{pro} \label{torsion fiber}
\source{geometric-bogomolov-conjecture-arbitrary-characteristics}
For any $t\in A'(\OK)_\tor\setminus A'_{e+1}(\OK)$, every irreducible component of the fiber $f^{-1}(t)\subset (X_{r-1}\times X)_\OK$ has canonical height 0 in $(A\times A)_\OK$.  
\end{pro}
\begin{proof}
Denote by $B$ the preimage of $A'$ under the summation map $A\times A\to A$. 
Since the preimage of $0$ under $A\times A\to A$ is geometrically integral, $B$ is also geometrically integral, and thus an abelian subvariety of $A\times A$ over $K$. We have the following commutative diagram.
$$
\xymatrix{
&\quad X_{r-1}\times X \quad  \ar[r] \ar[rd]^f  & \quad B \quad  \ar[d]^h \ar[r] &
\quad A\times A \quad \ar[d]
\\
&
&
\quad A' \quad \ar[r]
&
\quad A \quad
}
$$

Let $S/k$ with $K=k(S)$ be as in assumption (b). 
Denote by $\pi:\sA\to S$ the unique abelian scheme extending the abelian variety $A\to \Spec K$, which exists by assumption (d). 
Then $\sA\times_S\sA \to S$ is the unique abelian scheme extending $A\times A\to \Spec K$.
Denote by $\sX$ the Zariski closure of $X$ in $A$. 
Denote by $\sA'$ (resp. $\sB$) the Zariski closure of $A'$ in $\sA$ (resp. $B$ in $\sA\times_S\sA$). Both $\sA'$ and $\sB$ are abelian schemes over $S$. 
Denote by $\sY$ is the Zariski closure of $Y=X_{r-1}\times X$ in $\sA\times_S \sA$. 
This gives an integral version of the above diagram.
$$
\xymatrix{
\quad f^{-1}(\sT) \quad \ar[r] \ar[rd]
&\quad \sY \quad  \ar[r] \ar[rd]^f  & \quad \sB \quad  \ar[d]^h \ar[r] &
\quad \sA\times_S\sA \quad \ar[d]
\\
&
\quad \sT \quad \ar[r]
&
\quad \sA' \quad \ar[r]
&
\quad \sA \quad
}
$$

Let $\sT$ be a torsion multi-section of $\sA'\to S$ of order non-divisible by $\Char K$.
Then $\sT\to S$ is finite and \'etale, and $\sT$ is a smooth projective curve over $k$.
Assume that $T=\sT_K$ is not contained in $A'_{e+1}.$
Apply Proposition \ref{prostrleqtot} to the triangle of the second diagram. 
We obtain that 
$$\sum_{i=1}^n m_i[\sZ_i]\leq h^*[\sT]\cdot [\sY]$$
in $\CH_{e+1}(\sB)_\Q$. 
Here $\sZ_1, \dots, \sZ_n$ are irreducible components of $f^{-1}(\sT)$ satisfying 
$f(\sZ_i)=\sT$, and $m_i$ is the multiplicity of $\sZ_i$ in $f^{-1}(\sT).$ 

Let $\sL_{\sA'}$ (resp. $\sL_\sB$) be a symmetric, relatively ample and rigidified line bundle over $\sA'$ (resp. $\sB$).
By Proposition \ref{numerical equivalent}, there is a numerical equivalence 
$$[\sT]\equiv
a\, [\sL_{\sA'}]^{\dim A'}$$ 
in $\CH_{1}(\sA')_\Q$ for some $a>0$.

By these two results, we have
$$\sum_{i=1}^n m_i[\sZ_i]\cdot [\sL_\sB]^{e+1}
\leq a\,  h^*[\sT]\cdot [\sY] \cdot [\sL_\sB]^{e+1}
=a\, [h^*\sL_{\sA'}]^{\dim A'} \cdot [\sL_\sB]^{e+1}\cdot [\sY] .$$
By Lemma \ref{line bundle}(4), there is a constant $b>0$ such that $b\, \sL_\sB-h^*\sL_{\sA'}$ is a nef line bundle over $\sB$.
As a consequence, 
$$[h^*\sL_{\sA'}]^{\dim A'} \cdot [\sL_\sB]^{e+1} \cdot [\sY]
\leq 
b^{\dim A'}\,  [\sL_\sB]^{\dim \sY} \cdot [\sY].
$$

As in the proof of Lemma \ref{addition},  ``small points'' of $X$ transfer to ``small points'' of $Y=X_{r-1}\times X$.
Then $Y$ has a dense set of small points in $B$.
By Lemma \ref{fundamental inequality}, the height
$h_{\sL_\sB}(Y)=0$. 
Then we have $ [\sY] \cdot [\sL_\sB]^{\dim \sY}=0$. 
This forces $[\sZ_i]\cdot [\sL_\sB]^{e+1}$, and thus $h_{\sL_\sB}(Z_i)=0$.
Here $Z_i=\sZ_{i,K}$ for $i=1, \dots, n$ are exactly the irreducible components of $f^{-1}(\sT_K)$. 
This finishes the proof.
\end{proof}




\subsection{Lowering the dimension}

Now we prove that $X$ is torsion by induction on $\dim X$.
If $\dim X=0$, we must have $\hat h(X)=0$, and then $X$ is a torsion point of $A$ by the Northcott theorem (cf. \cite[Theorem 9.15]{Conrad}).  

If $\dim X>0$, consider the morphism $f:X_{r-1}\times X\to A'$ obtained in assumption (e). 
By Proposition \ref{torsion fiber}, for any $t\in A'(\OK)_\tor\setminus A'_{e+1}(\OK)$, every irreducible component of the fiber $f^{-1}(t)\subset (X_{r-1}\times X)_\OK$ has canonical height 0 in $A\times A$.  
Note that every irreducible component of $f^{-1}(t)$ has dimension $e<\dim X$.
By induction, it is torsion in $A\times A$. 
Therefore, $f^{-1}(t)$ contains a Zariski dense set of torsion points of $A\times A$. 
As $A'(\OK)_\tor\setminus A'_{e+1}(\OK)$ is Zariski dense in $A'$,
$X_{r-1}\times X$ contains a Zariski dense set of torsion points of $A\times A$.  

By the Manin--Mumford conjecture (cf. Theorem \ref{MM}), 
$X_{r-1}\times X$ is a torsion subvariety of $A\times A$.  
Then $X$ is a torsion subvariety of $A$, since it is the image of the composition of $X_{r-1}\times X\to A\times A$ with $p_2:A\times A\to A$. 
This finishes the proof of Theorem \ref{thmmain}. 








%\newpage

\bibliography{dd}
